\documentclass[11pt]{article}
\usepackage[margin=1in]{geometry}
\usepackage{amsmath,amssymb,amsthm,mathtools,booktabs,microtype}
\usepackage[hidelinks]{hyperref}
\hypersetup{pdftitle={Sharp reference-assisted contraction for the minimal quantum Fisher metric},pdfauthor={Lluis Eriksson},pdfsubject={Complete equal-reference contraction, random-unitary channels, and Petz divergences},pdfkeywords={quantum information, contraction coefficient, SLD, Petz divergence, depolarizing channel}}
\newcommand{\Tr}{\operatorname{Tr}}
\newcommand{\id}{\operatorname{id}}
\newcommand{\Ad}{\operatorname{Ad}}
\newcommand{\E}{\mathcal E}
\newcommand{\N}{\mathcal N}
\newcommand{\ip}[2]{\langle #1,#2\rangle}
\newcommand{\Dtri}{D_{\triangle}}
\newcommand{\gmin}{g^{\mathrm{SLD}}}
\newtheorem{theorem}{Theorem}
\newtheorem{proposition}[theorem]{Proposition}
\newtheorem{corollary}[theorem]{Corollary}
\newtheorem{lemma}[theorem]{Lemma}
\theoremstyle{remark}
\newtheorem{remark}[theorem]{Remark}
\title{Sharp reference-assisted contraction for the minimal quantum Fisher metric}
\author{Lluis Eriksson}
\date{27 October 2026}
\begin{document}
\maketitle
\begin{abstract}
We derive a harmonic-weight contraction bound for random-unitary channels on pairs with the same uniform-twirl image. It applies to every operator-convex Petz divergence on faithful states and is uniform over finite quantum reference systems. For a qubit depolarizing channel with retained amplitude \(0<\lambda<1\), we identify the exact complete equal-reference contraction coefficient of both the minimal quantum Fisher metric and the triangular Petz divergence:
\(3\lambda^2/(1+2\lambda)\).
The lower bound uses an explicit qubit--qutrit ladder state, a partial-trace-zero tangent, and an ordered faithful regularization. Thus the general Petz upper bound is sharp as a bound uniform over that family. At \(\lambda=1/2\) it equals \(3/8\). This does not determine the relative-entropy or mutual-information coefficient. We include a finite rational triangular-divergence witness above \(1/3\), explain a limitation of scalar-centering proofs, and preserve failed auxiliary variance and numerical routes.
\end{abstract}

\section{Quantities and scope}
Let \(A\) and \(R\) be finite-dimensional systems. All local maps below are tensored with \(\id_R\), whether or not this is written. A state is faithful if it is positive definite. We use the Hilbert--Schmidt inner product
\(\ip{X}{Y}=\Tr X^\dagger Y\) on the \emph{complex} Hilbert space of matrices. Define
\(L_a(X)=aX\) and \(R_b(X)=Xb\).
For faithful \(a,b\), these are commuting positive invertible superoperators. This statement concerns positivity on Hilbert--Schmidt space, not preservation of the cone of positive matrices.

For a faithful state \(\sigma\) and a Hermitian tangent \(Z\), set
\begin{equation}\label{eq:sld}
 \gmin_\sigma(Z)=
 \ip{Z}{\bigl[(L_\sigma+R_\sigma)/2\bigr]^{-1}Z}.
\end{equation}
This is the minimal monotone, or symmetric logarithmic derivative (SLD), quantum Fisher quadratic form. Its associated finite chi-square divergence is
\begin{equation}\label{eq:chi}
 \chi^2_{\min}(\rho\Vert\sigma)=\gmin_\sigma(\rho-\sigma).
\end{equation}
For an operator-convex \(f:(0,\infty)\to\mathbb R\), \(f(1)=0\), the Petz divergence is
\begin{equation}\label{eq:petz}
 D_f(\rho\Vert\sigma)=
 \ip{\sigma^{1/2}}{f(L_\rho R_\sigma^{-1})\sigma^{1/2}}.
\end{equation}
We will also use the distinct finite divergence
\begin{equation}\label{eq:triangle}
 f_\triangle(x)=\frac{(x-1)^2}{x+1},\qquad
 \Dtri(\rho\Vert\sigma)=D_{f_\triangle}(\rho\Vert\sigma).
\end{equation}
Although \eqref{eq:chi} and \eqref{eq:triangle} have proportional local forms, they are not the same finite divergence. Nor do we identify either with the Bures distance or its geodesic contraction.

For a local channel \(\N\), define the complete equal-reference coefficient
\begin{equation}\label{eq:complete}
 \eta_f^c(\N)=\sup_{\substack{|R|<\infty,\ \rho,\sigma>0\\
 \rho\ne\sigma,\ \Tr_A\rho=\Tr_A\sigma}}
 \frac{D_f(\N(\rho)\Vert\N(\sigma))}
      {D_f(\rho\Vert\sigma)}.
\end{equation}
Define its SLD analogue by optimizing
\(\gmin_{\N(\sigma)}(\N(Z))/\gmin_\sigma(Z)\)
over faithful \(\sigma\), nonzero Hermitian \(Z\) with \(\Tr_AZ=0\), and all finite \(R\).
It is also the coefficient of \eqref{eq:chi}: every such tangent is realized by a sufficiently small faithful pair.

These restrictions matter. Actual mutual information compares a state with the product of \emph{its own} marginals. Conditional coefficients can instead subtract a divergence between unequal reference marginals. Fixed-point entropy decay compares a state with a conditional expectation. We do not replace one optimization by another.

\section{Known representation and a cancellation bound}
The positive resolvent representation of operator-convex functions, and the resulting representation of Petz divergences, are established facts; see Hiai--Ruskai \cite{HR}, Proposition 2.1 and equations (12)--(13). After removing an affine term, any nonaffine operator-convex \(f\) as above has
\begin{equation}\label{eq:representation}
 f(x)=f'(1)(x-1)+c(x-1)^2+
 \int_{[0,\infty)}\frac{(x-1)^2}{x+t}\,d\mu(t),
\end{equation}
where \(c\ge0\), \(\mu\) is positive and
\(\int(1+t)^{-1}\,d\mu(t)<\infty\).
Writing \(Z=\rho-\sigma\), for states of the same trace this gives
\begin{equation}\label{eq:resolvent}
 D_f(\rho\Vert\sigma)=
 c\,\Tr Z^2\sigma^{-1}+
 \int_{[0,\infty)}
 \ip{Z}{(L_\rho+tR_\sigma)^{-1}Z}\,d\mu(t).
\end{equation}
We do not claim this representation as a contribution.

\begin{lemma}\label{lem:inverse}
For positive definite \(Q\), the inverse quadratic form
\(F(Q,Z)=\ip{Z}{Q^{-1}Z}\) is jointly convex in \(Q,Z\), for arbitrary complex vectors \(Z\).
\end{lemma}
\begin{proof}
Use the variational identity
\[
 F(Q,Z)=\sup_Y\{2\operatorname{Re}\ip{Z}{Y}-\ip{Y}{QY}\}.
\]
Each expression inside the supremum is affine in \(Q,Z\).
The same identity also shows that for any trial \(Y\), with \(e=Z-QY\),
\begin{equation}\label{eq:residual}
 F(Q,Z)=2\operatorname{Re}\ip{Z}{Y}-\ip{Y}{QY}
       +\ip{e}{Q^{-1}e}.
\end{equation}
\end{proof}

\begin{theorem}[Harmonic cancellation]\label{thm:harmonic}
Let \(U_1,\ldots,U_m\) be local unitaries, \(m\ge2\), and let
\[
 \E=\frac1m\sum_{j=1}^m\Ad_{U_j},\qquad
 \N_w=\sum_{j=1}^m w_j\Ad_{U_j},\qquad
 w_j>0,\quad \sum_jw_j=1.
\]
For faithful states \(\rho,\sigma\) satisfying \(\E(\rho)=\E(\sigma)\),
and every nonaffine operator-convex \(f\) with \(f(1)=0\),
\begin{equation}\label{eq:harmonic}
 D_f(\N_w(\rho)\Vert\N_w(\sigma))
 \le h(w)D_f(\rho\Vert\sigma),\qquad
 h(w)=1-\frac{m^2}{\sum_jw_j^{-1}}.
\end{equation}
The same \(h(w)\) applies for every finite \(R\).
It also bounds the SLD metric ratio for tangents \(Z\) with \(\E(Z)=0\).
\end{theorem}
\begin{proof}
For a resolvent term in \eqref{eq:resolvent}, let
\(Q=L_\rho+tR_\sigma\), \(Z=\rho-\sigma\), and write
\(Q_j=\Ad_{U_j}Q\Ad_{U_j}^\dagger\), \(Z_j=U_jZU_j^\dagger\).
The constraint gives \(\sum_jZ_j=0\).
For every real \(a\),
\[
 \N_w(Z)=\sum_j(w_j-a)Z_j
        =\sum_jw_j\frac{w_j-a}{w_j}Z_j.
\]
Since the output superoperator is \(\sum_jw_jQ_j\), Lemma \ref{lem:inverse}, quadratic homogeneity and unitary covariance give
\[
 F\left(\sum_jw_jQ_j,\N_w(Z)\right)
 \le \sum_j\frac{(w_j-a)^2}{w_j}F(Q,Z).
\]
The scalar factor is
\(1-2ma+a^2\sum_jw_j^{-1}\); its minimum occurs at
\(a=m/\sum_jw_j^{-1}\) and equals \(h(w)\).
The convex weights here are the positive \(w_j\). Negative tangent multipliers do not constitute negative convex weights on states.

For the quadratic endpoint in \eqref{eq:resolvent}, use \(Q=R_\sigma\); its inverse quadratic form is \(\Tr Z^2\sigma^{-1}\).
Integrate the nonnegative resolvent inequalities. Affine terms vanish by equal traces, and the positive representation justifies this passage.
For the SLD metric, use \(Q=(L_\sigma+R_\sigma)/2\) directly.
No argument depends on a bound on \(|R|\).
\end{proof}

Cauchy--Schwarz gives \(\sum_jw_j^{-1}\ge m^2\), so \(0\le h(w)<1\).
A useful exact control is an involution \(U=U^\dagger\), \(U^2=I\), with both eigenspaces nonzero. For \(\E=(\id+\Ad_U)/2\) and
\(\N=w\id+(1-w)\Ad_U\), the bound is \((2w-1)^2\).
It is sharp for each of the Petz generators under consideration: take
\(\sigma=I/n\), a nonzero Hermitian \(Z\) connecting the two eigenspaces, and
\(\rho_\epsilon=\sigma+\epsilon Z\).
Then \(\N(Z)=(2w-1)Z\), and the scalar second-order expansion of \eqref{eq:petz} gives the stated limiting ratio. At \(w=1/2\), the two outputs coincide.

\begin{remark}[Relative entropy at the boundary]\label{rem:kl}
For \(f(x)=x\log x\), \eqref{eq:harmonic} also holds for nonfaithful pairs with finite relative entropy. Mix both states with the same maximally mixed state. The equal-\(\E\) constraint is preserved. Joint convexity and lower semicontinuity imply convergence of both finite divergences, and the faithful inequality passes to the limit before division. We do not assert a boundary extension for every \(f\) without specifying its boundary conventions.
\end{remark}

\section{Depolarization and the exact minimal-metric coefficient}
For a \(d\)-dimensional unitary error basis, \(m=d^2\), the uniform twirl is the local replacer
\(\E_A(X)=I_A/d\otimes\Tr_AX\).
With
\(w_0=[1+(m-1)\lambda]/m\) and
\(w_j=(1-\lambda)/m\) for \(j\ne0\),
Theorem \ref{thm:harmonic} gives
\begin{equation}\label{eq:dim}
 h_d(\lambda)=\frac{(d^2-1)\lambda^2}{1+(d^2-2)\lambda}.
\end{equation}
Equal twirl images are precisely equal reference marginals.
For qubits, \eqref{eq:dim} is sharp for the minimal metric and \eqref{eq:triangle}.

\begin{theorem}[Exact complete coefficients for a qubit]\label{thm:exact}
Let \(0<\lambda<1\) and
\(\N_\lambda(X)=\lambda X+(1-\lambda)I_A/2\otimes\Tr_AX\), with \(|A|=2\).
Then
\begin{equation}\label{eq:exact}
 \eta_{\mathrm{SLD}}^c(\N_\lambda)
 =\eta_{\chi^2_{\min}}^c(\N_\lambda)
 =\eta_{\triangle}^c(\N_\lambda)
 =\frac{3\lambda^2}{1+2\lambda}.
\end{equation}
A reference of dimension three suffices to approach these suprema.
No assertion is made that three is the smallest sufficient dimension.
\end{theorem}
\begin{proof}
The upper bounds follow from Theorem \ref{thm:harmonic} with the four Pauli conjugations and their depolarizing weights. It remains to prove the lower bounds.

In the ordered basis
\((|00\rangle,|01\rangle,|02\rangle,|10\rangle,|11\rangle,|12\rangle)\),
fix \(e>0\) and define
\[
 a=(1,0,0,0,e,0)^T,\quad
 b=(0,e,0,0,0,1)^T,\quad
 \sigma_e=\frac{aa^\dagger+bb^\dagger}{2(1+e^2)}.
\]
The vectors are orthogonal and have the same norm, so
\(P_e=2\sigma_e\) is a rank-two projection.
Let \(Z_e\) and \(J\) be
\[
 Z_e=\begin{pmatrix}0&-1\\-1&-2e\end{pmatrix},
 \qquad J=\begin{pmatrix}0&1\\1&0\end{pmatrix}.
\]
Define the Hermitian \(6\times6\) matrix \(\Delta_e\) by its two nonzero blocks:
\(Z_e\) on indices \((0,4)\) and
\(-JZ_eJ=\bigl(\begin{smallmatrix}2e&1\\1&0\end{smallmatrix}\bigr)\)
on \((1,5)\). All other entries vanish.
Direct multiplication gives
\begin{align}
 \Tr_A\Delta_e&=0,&
 (I-P_e)\Delta_e(I-P_e)&=0,\label{eq:kernel}\\
 \Tr\Delta_e^2&=4+8e^2,&
 \Tr(P_e\Delta_eP_e)^2&=8e^2.\label{eq:traces}
\end{align}
The SLD form at this nonfaithful base is understood by retaining the support--support and support--kernel terms; the kernel--kernel terms vanish by \eqref{eq:kernel}. The two positive eigenvalues are \(1/2\), so \eqref{eq:traces} yields
\begin{equation}\label{eq:input}
 \gmin_{\sigma_e}(\Delta_e)
 =4\Tr\Delta_e^2-2\Tr(P_e\Delta_eP_e)^2
 =16(1+e^2).
\end{equation}

The reference marginal is
\[
 \sigma_{e,R}=\frac{\operatorname{diag}(1,2e^2,1)}{2(1+e^2)}.
\]
On indices \((0,4)\), the output base is
\begin{equation}\label{eq:block}
 W_e=\frac1{4(1+e^2)}
 \begin{pmatrix}1+\lambda&2\lambda e\\2\lambda e&2e^2\end{pmatrix}.
\end{equation}
On \((1,5)\) it is \(JW_eJ\).
The other two singleton directions have zero tangent.
For \(e>0\) and \(0<\lambda<1\), the output is faithful:
the block determinant is positive because
\(1+\lambda-2\lambda^2=(1-\lambda)(1+2\lambda)>0\),
and the singleton entries are positive.

Solving the two-dimensional Sylvester equation
\(W_eY+YW_e=2Z_e\), and using the two mirrored blocks, gives
\begin{equation}\label{eq:output}
 \gmin_{\N_\lambda(\sigma_e)}(\lambda\Delta_e)
 =
 \frac{16\lambda^2(1+e^2)
 (4e^2\lambda+2e^2+3\lambda+3)}
 {(2\lambda+1)(2e^2+\lambda+1)}.
\end{equation}
This is a \(2\times2\) linear calculation, not an optimization claim.
Dividing by \eqref{eq:input} gives
\begin{equation}\label{eq:ratio}
 R_e(\lambda)=
 \frac{\lambda^2(4e^2\lambda+2e^2+3\lambda+3)}
 {(2\lambda+1)(2e^2+\lambda+1)}.
\end{equation}
In particular,
\begin{equation}\label{eq:gap}
 \frac{3\lambda^2}{1+2\lambda}-R_e(\lambda)
 =
 \frac{4e^2\lambda^2(1-\lambda)}
 {(2\lambda+1)(2e^2+\lambda+1)}
 \longrightarrow0 \quad(e\downarrow0).
\end{equation}

We now realize the limits through faithful inputs. For fixed \(e>0\), set
\(\sigma_{e,\xi}=(1-\xi)\sigma_e+\xi I_6/6\), \(0<\xi<1\).
Its eigenspaces are those of \(\sigma_e\).
Equation \eqref{eq:kernel} implies that its input SLD form converges to
\eqref{eq:input} as \(\xi\downarrow0\); there are no diverging kernel--kernel contributions.
The output tends to the faithful \(\N_\lambda(\sigma_e)\), so its SLD form converges by finite-dimensional continuity.
Consequently faithful tangent ratios approach \(R_e(\lambda)\).
Sending \(e\downarrow0\) afterwards proves the SLD lower bound.
The equality with the complete chi-square coefficient follows from the exact quadratic dependence in \eqref{eq:chi} and sufficiently small faithful perturbations.

For the triangular Petz divergence, at each fixed faithful \(\sigma_{e,\xi}\) choose
\(\rho_t=\sigma_{e,\xi}+t\Delta_e\).
For sufficiently small nonzero real \(t\), this is a faithful density matrix with the same reference marginal.
Its resolvent representation gives
\[
 \lim_{t\to0}\frac{\Dtri(\rho_t\Vert\sigma_{e,\xi})}{t^2}
 =
 \ip{\Delta_e}{(L_{\sigma_{e,\xi}}+R_{\sigma_{e,\xi}})^{-1}\Delta_e}
 =\frac12\gmin_{\sigma_{e,\xi}}(\Delta_e),
\]
and the corresponding output limit is analogous.
Thus take \(t\to0\), then \(\xi\downarrow0\), and finally \(e\downarrow0\).
This supplies the lower supremum over faithful finite pairs in \eqref{eq:complete}.
There is no exchange of these limits and no assumed uniform continuity at a singular reference.
\end{proof}

\begin{corollary}\label{cor:uniform}
For the qubit depolarizing channel, \(3\lambda^2/(1+2\lambda)\) is the smallest bound valid uniformly over all nonaffine operator-convex Petz generators and faithful equal-reference pairs.
\end{corollary}
\begin{proof}
Theorem \ref{thm:harmonic} gives the uniform upper; the member \(f_\triangle\) attains the supremum in Theorem \ref{thm:exact}.
This does not assert that every generator has that coefficient.
\end{proof}

At \(\lambda=1/2\), \eqref{eq:exact} equals \(3/8\).
The ordinary unital-qubit contraction coefficient is \(\lambda^2\), by Hiai--Ruskai \cite{HR}, Theorem 6.1. At \(\lambda=1/2\), the ordinary value is \(1/4\), compared with the complete equal-reference value \(3/8\).
The stabilization considered here retains equal reference marginals and optimizes over their quantum correlations.

\section{A limitation of changing the unitary decomposition}
A narrow method limitation helps distinguish the remaining relative-entropy problem from a representation search.

\begin{proposition}\label{prop:barrier}
Fix \(0<\lambda<1\) and let \(E(X)=\Tr(X)I/d\).
Represent \(E=\sum_{j=1}^nq_j\Ad_{U_j}\) as a finite convex combination, with \(q_j\ge0\) and \(\sum_jq_j=1\); discard zero-weight atoms. Consolidate all identity conjugations into an atom of total weight \(r\); if absent, take \(r=0\).
Use the induced common representation of
\(N=\lambda\id+(1-\lambda)E\).
Among all such representations, a single scalar cancellation of the form
\(w_j-aq_j\), optimized by inverse-quadratic convexity, cannot give an upper smaller than \eqref{eq:dim}.
An error basis attains that method-family minimum.
\end{proposition}
\begin{proof}
The scalar minimum is
\(1-1/S\), \(S=\sum_jq_j^2/w_j\), because
\(\sum_j(w_j-aq_j)^2/w_j=1-2a+a^2S\).
The induced weights are \(w_0=\lambda+(1-\lambda)r\) and
\(w_j=(1-\lambda)q_j\) otherwise. Hence
\[
 S=\frac{r^2}{\lambda+(1-\lambda)r}+\frac{1-r}{1-\lambda},
 \qquad
 h(r)=\frac{\lambda^2(1-r)}{\lambda+(1-2\lambda)r}.
\]
Using normalized Choi states, \(J(E)=I_{d^2}/d^2\) and
\(J(\id)=|\Omega\rangle\langle\Omega|\).
Positivity of the other Choi terms gives
\(J(E)-r|\Omega\rangle\langle\Omega|\succeq0\), so \(r\le1/d^2\).
On this interval,
\[
 h'(r)=-\frac{\lambda^2(1-\lambda)}
 {[\lambda+(1-2\lambda)r]^2}<0.
\]
The smallest possible scalar bound therefore occurs at \(r=1/d^2\) and equals \eqref{eq:dim}, attained by a unitary error basis.
\end{proof}
This concerns the specified common-representation scalar method. It does not exclude other representations, state-dependent cancellations, stronger operator inequalities, or a smaller true coefficient for relative entropy.

\section{Finite certificates, failures, and unresolved questions}
The accompanying evidence distinguishes analytical arguments from finite diagnostics.
The exact symbolic checks verify the projector and partial-trace identities in \eqref{eq:kernel}--\eqref{eq:traces}, the channel block \eqref{eq:block}, the Sylvester equation, and the rational ratio and gap.
Nine faithful finite pairs were also evaluated at 75 decimal digits. They are controls on the construction, not the proof of its universal upper or sharp supremum.

A separate finite rational certificate at \(\lambda=1/2\) gives faithful qubit--qutrit states with the same reference marginal and
\begin{equation}\label{eq:finite}
 \frac{\Dtri(\N(\rho)\Vert\N(\sigma))}
 {\Dtri(\rho\Vert\sigma)}>0.348913969797>\frac13.
\end{equation}
The certificate stores the matrices and rational trial Sylvester solutions.
If \(Q\succeq cI\), \eqref{eq:residual} bounds the remainder by
\(\|e\|_2^2/c\).
Exact rational input/output intervals then certify \eqref{eq:finite}, without treating a floating-point eigenvalue as a certificate.

We also preserve a failed route through conditional variance. For
\[
 V_\sigma(H)=\inf_{B=B^\dagger}\Tr\sigma(H-I_A\otimes B)^2,
 \qquad \Tr_AH=0,
\]
a proposed sufficient inequality was
\(V_\sigma(H)\le2\Tr[E_A(\sigma)H^2]\).
A rational qubit--qutrit certificate instead proves a ratio greater than \(2.99311283224\).
The minimizer solves a reference Sylvester equation, and the same residual method certifies the violation.
This refutes the auxiliary inequality; it does not directly refute an SLD channel coefficient or the relative-entropy conjecture.

Exploration included fixed-budget local optimizations, some terminating at their limits or abnormally, and unsuccessful simpler families.
These outcomes are retained. The reversed ladder and the proof above replace reliance on optimized values.
No exhaustive computational search or independently certified computational campaign is claimed.

The relative-entropy complete coefficient at \(\lambda=1/2\) is not determined here.
In particular, the triangular example cannot be transferred to BKM by comparing the sizes of their metrics: ordering metrics does not order their contraction ratios.
Our results identify an obstruction to a bound \emph{uniform over all Petz generators}, while leaving a generator-specific relative-entropy improvement open.
No capacity, tensorization, optimal statistical protocol, geodesic-distance contraction or attained singular maximum is asserted.


\paragraph{Evidence availability and replay.}
A \href{https://github.com/lluiseriksson/physmath-knowledge-tree/tree/f0b3642b196b9de6d51b31cca49a69818522ee91/integrations/research-radar/evidence/2026-10-27-pauli-audit}{version-pinned public supplement}
contains the scripts, rational certificates, complete component report and retained objections.
Its \texttt{MANIFEST1001.json} binds the LF transport files to their sealed-original hashes.
Copy the files to a temporary directory, install the pinned \texttt{requirements.txt}, and run
\texttt{python symbolic993.py}, \texttt{python verify993.py}, and
\texttt{python verify-certificates991.py}. These programs write their reports beside the scripts.
The later, explicitly speculative Pauli extension in the same research collection is outside the claims and assessment of this manuscript.

\section{Relation to prior work}
Hiai--Ruskai \cite{HR} develop contraction comparisons, monotone metrics and ordinary qubit formulas. Their resolvent representation is the starting tool here.
Cao--Lu \cite{CL} study quantum chi-square SDPIs, tensorization with specified product reference states and ordinary qubit examples. Those hypotheses are distinct from optimization over arbitrary quantum references in \eqref{eq:complete}.
George et al. \cite{Doeblin} formulate equal-reference coefficients and bound complete trace-distance contraction through quantum Doeblin coefficients. George--Tomamichel \cite{GT} treat fixed-base and ordinary optimized chi-square coefficients and associated spectral algorithms. Their fixed-base optimization does not itself determine the supremum over correlated bases and all reference dimensions under \(\Tr_A Z=0\). The unrestricted ordinary coefficient of \(\N\otimes\id_R\) can instead equal one because a tangent may reside solely in \(R\); it is not the coefficient in \eqref{eq:complete}.
Belzig et al. \cite{Reverse} distinguish ordinary channel comparisons from quantum-reference mutual-information questions.
Hirche et al. \cite{Conditional} distinguish mutual-information, complete equal-reference and conditional coefficients; these distinctions prevent interpreting our triangular result as a relative-entropy or mutual-information solution.
Gao--Zhao \cite{GZ} concern entropy contraction toward a conditional expectation, another important nearby optimization that is not identified with \eqref{eq:complete}.

These references are included to delimit the quantities and established machinery, not to suggest that the surrounding literature was exhaustively audited.
The bibliographic search was directed. Certified worldwide originality, independent human review and formal proof verification are not claimed.

\section{Model assessments and review context}
The named author directed substantial OpenAI Codex assistance in derivation, writing and numerical work.
Technical audit 988 and full-draft review 998 were requested with GPT-6 Astra at xhigh effort in reused conversational contexts and a shared filesystem. Requested configuration does not attest effective runtime identity, isolation or independence. Review 998 evaluated the earlier exact seven-page draft; the present corrections address its supplement, bibliography and hypothesis-precision objections. Its score is not automatically assigned to this corrected PDF.
Component checks do not constitute a full assessment of this exact manuscript.
Any actual reports must identify their exact input hashes, retained objections, context reuse and scope.
No score from an earlier manuscript is transferred to this text; no predicted assessment score is claimed.
The author remains responsible for the claims, attribution and publication decision.
This section distinguishes AI assessment, arithmetic replay, formal verification, human peer review and editorial acceptance.

\small
\begin{thebibliography}{9}
\bibitem{HR}
F. Hiai and M. B. Ruskai,
\emph{Contraction coefficients for noisy quantum channels},
arXiv:1508.03551v1 (2015).
\url{https://arxiv.org/abs/1508.03551v1}.

\bibitem{CL}
Y. Cao and J. Lu,
\emph{Tensorization of the strong data processing inequality for quantum chi-square divergences},
Quantum \textbf{3}, 199 (2019).
\url{https://doi.org/10.22331/q-2019-10-28-199}.

\bibitem{Doeblin}
I. George, C. Hirche, T. Nuradha and M. M. Wilde,
\emph{Quantum Doeblin Coefficients: Interpretations and Applications},
Quantum \textbf{10}, 2115 (2026), arXiv:2503.22823v3.
\url{https://arxiv.org/abs/2503.22823v3}.

\bibitem{Reverse}
P. Belzig, L. Gao, G. Smith and P. Wu,
\emph{Reverse-type Data Processing Inequality},
arXiv:2411.19890v3 (2025).
\url{https://arxiv.org/abs/2411.19890v3}.

\bibitem{Conditional}
C. Hirche, I. George, T. Nuradha and M. M. Wilde,
\emph{Conditional contraction coefficients and their applications to quantum networks},
arXiv:2608.27171v1 (2026).
\url{https://arxiv.org/abs/2608.27171v1}.

\bibitem{GZ}
L. Gao and L. Zhao,
\emph{Tight Entropy Contraction of Generalized Quantum Depolarization},
arXiv:2608.19685v1 (2026).
\url{https://arxiv.org/abs/2608.19685v1}.
\bibitem{GT}
I. George and M. Tomamichel,
\emph{A Unified Approach to Quantum Contraction and Correlation Coefficients},
arXiv:2505.15281v3 (2026).
\url{https://arxiv.org/abs/2505.15281v3}.

\end{thebibliography}
\end{document}
